\section*{Глава 2. Проектирование веб-приложения}
\addcontentsline{toc}{section}{Глава 2. Проектирование веб-приложения}
\stepcounter{section}


\subsection{Анализ целевой аудитории}

Разрабатываемое веб-приложение ориентировано на пользователей, которые хотят 
сохранять и систематизировать цитаты из фильмов. В ходе анализа были выделены два 
основных типа пользователей: обычные пользователи и администраторы.

Обычные пользователи — это основная аудитория приложения. К ней относятся люди, 
которые смотрят фильмы, запоминают яркие фразы и хотели бы сохранять их в одном месте. 
Существующие способы (заметки в телефоне, файлы на компьютере, бумажные блокноты) 
не обеспечивают удобной систематизации и аналитики.

Основные потребности обычного пользователя:
\begin{itemize}
    \item быстрое добавление цитаты с указанием фильма, персонажа и таймкода;
    \item просмотр, редактирование и удаление своих цитат;
    \item поиск и фильтрация по фильмам или персонажам;
    \item визуализация личной статистики (например, какие фильмы чаще всего цитируются);
    \item импорт цитат из файла (CSV/JSON) и экспорт коллекции для бэкапа.
\end{itemize}

На основе этих потребностей формируется функциональность приложения: личный кабинет, 
CRUD-операции с цитатами, модуль аналитики, импорт и экспорт данных.

Администратор — это пользователь с расширенными правами, который поддерживает каталог 
фильмов в аккуратном состоянии.

Задачи администратора:
\begin{itemize}
    \item редактирование информации о любых фильмах (исправление названия, года, жанра);
    \item объединение дублирующихся записей об одном и том же фильме;
    \item управление справочником жанров (добавление, переименование).
\end{itemize}

Таким образом, наличие двух ролей позволяет не только обеспечить безопасность данных, но и 
поддерживать высокое качество общего каталога фильмов.

\subsection{Описание функциональности приложения}

На основе анализа целевой аудитории был сформулирован основной сценарий использования 
системы в виде User Story: «Я, как обычный пользователь, хочу сохранять цитаты из фильмов 
с указанием фильма, персонажа и таймкода, чтобы в любой момент вернуться к любимым фразам 
и видеть статистику по своей коллекции». Данный сценарий определяет ключевые требования 
к функциональности приложения.

На рисунке~\ref{fig:usecase} представлена диаграмма вариантов использования (UseCase), 
отображающая действия, которые могут выполнять пользователи в системе.

\begin{figure}[ht!]
    \centering
    \includegraphics[width=0.9\textwidth]{usecase.png}
    \caption{Диаграмма вариантов использования}
    \label{fig:usecase}
\end{figure}

\textbf{Функциональность обычного пользователя:}
\begin{enumerate}
    \item Регистрация и вход в систему.
    \item Просмотр списка своих цитат.
    \item Добавление новой цитаты (с выбором или добавлением фильма).
    \item Редактирование и удаление своих цитат.
    \item Просмотр личной аналитики (топ фильмов, распределение по жанрам, динамика по месяцам).
    \item Импорт цитат из файлов (CSV/JSON).
    \item Экспорт коллекции в файл (CSV/JSON).
\end{enumerate}

\textbf{Дополнительная функциональность администратора:}
\begin{enumerate}
    \item Редактирование и удаление любых фильмов.
    \item Управление жанрами.
    \item Просмотр общей аналитики по системе.
\end{enumerate}


\subsection{Проектирование модели данных}

База данных состоит из следующих основных сущностей: пользователи, жанры, фильмы, цитаты.

\textbf{Таблица users (пользователи):}
\begin{itemize}
    \item id (SERIAL, PRIMARY KEY)
    \item login (VARCHAR(100), NOT NULL, UNIQUE)
    \item password\_hash (VARCHAR(256), NOT NULL)
    \item role (VARCHAR(20), NOT NULL, DEFAULT 'user')
    \item created\_at (TIMESTAMP, DEFAULT NOW())
\end{itemize}

\textbf{Таблица genres (жанры):}
\begin{itemize}
    \item id (SERIAL, PRIMARY KEY)
    \item name (VARCHAR(100), NOT NULL, UNIQUE)
\end{itemize}

\textbf{Таблица movies (фильмы):}
\begin{itemize}
    \item id (SERIAL, PRIMARY KEY)
    \item title (VARCHAR(255), NOT NULL)
    \item release\_year (INTEGER)
    \item director (VARCHAR(255))
    \item genre\_id (INTEGER, FOREIGN KEY REFERENCES genres)
    \item added\_by (INTEGER, FOREIGN KEY REFERENCES users)
    \item created\_at (TIMESTAMP, DEFAULT NOW())
\end{itemize}

\textbf{Таблица quotes (цитаты):}
\begin{itemize}
    \item id (SERIAL, PRIMARY KEY)
    \item text (TEXT, NOT NULL)
    \item character\_name (VARCHAR(255))
    \item timestamp (VARCHAR(20))
    \item movie\_id (INTEGER, FOREIGN KEY REFERENCES movies)
    \item user\_id (INTEGER, FOREIGN KEY REFERENCES users)
    \item created\_at (TIMESTAMP, DEFAULT NOW())
\end{itemize}

Связи между таблицами:
\begin{itemize}
    \item Один фильм может иметь много цитат (1 : N).
    \item Один пользователь может добавить много цитат (1 : N).
    \item Один жанр может относиться к многим фильмам (1 : N).
\end{itemize}

ER-диаграмма представлена на рисунке~\ref{fig:erd}.

\begin{figure}[ht!]
    \centering
    \includegraphics[width=0.8\textwidth]{erd.png}
    \caption{ER-диаграмма базы данных}
    \label{fig:erd}
\end{figure}

\subsection{Разработка макетов страниц}
\newcommand{\imgwidth}{12cm}

Были разработаны макеты следующих страниц: страница входа, страница со списком 
цитат, страница добавления цитаты, страница аналитики, страница администратора.

\textbf{Страница входа.} Пользователь вводит логин и пароль. При нажатии на 
ссылку «Зарегистрироваться» открывается форма регистрации.

\begin{figure}[ht!]
    \centering
    \includegraphics[width=\imgwidth]{login.png}
    \caption{Макет страницы входа}
    \label{fig:login}
\end{figure}

\textbf{Страница со списком цитат.} После входа пользователь видит свои цитаты 
в виде карточек. Над списком размещены кнопки «Добавить цитату», «Импорт», «Экспорт», «Аналитика».
Макет главной страницы приведён для обычного пользователя. Администратор видит ту же страницу, 
но с дополнительной кнопкой "Админ-панель" в верхней панели.

\begin{figure}[ht!]
    \centering
    \includegraphics[width=\imgwidth]{quotes_list.png}
    \caption{Макет страницы со списком цитат}
    \label{fig:quotes}
\end{figure}

\textbf{Страница добавления цитаты.} Пользователь выбирает фильм из выпадающего 
списка (или добавляет новый), вводит текст цитаты, персонажа и таймкод (необязательно).

\begin{figure}[ht!]
    \centering
    \includegraphics[width=\imgwidth]{add_quote.png}
    \caption{Макет страницы добавления цитаты}
    \label{fig:add_quote}
\end{figure}

\textbf{Страница аналитики.} На странице отображаются три графика: топ фильмов 
по цитатам (гистограмма), распределение по жанрам (круговая диаграмма), динамика 
по месяцам (линейный график).

\begin{figure}[ht!]
    \centering
    \includegraphics[width=\imgwidth]{analytics.png}
    \caption{Макет страницы аналитики}
    \label{fig:analytics}
\end{figure}

\textbf{Панель администратора (управление фильмами)} Для администратора предусмотрен отдельный 
интерфейс, доступный по кнопке «Админ-панель» в верхней панели навигации. 
Данная страница содержит две вкладки: «Фильмы» и «Жанры».

На вкладке «Фильмы» администратор может просматривать все фильмы, добавленные пользователями, 
редактировать их данные (название, режиссёр, год, жанр) или удалять ошибочные записи. 


\begin{figure}[ht!]
    \centering
    \includegraphics[width=\imgwidth]{admin_movies.png}
    \caption{Макет панели администратора (вкладка «Фильмы»)}
    \label{fig:admin_movies}
\end{figure}

На вкладке «Жанры» доступно управление справочником жанров: добавление новых, переименование 
и удаление.

\begin{figure}[ht!]
    \centering
    \includegraphics[width=\imgwidth]{admin_genres.png}
    \caption{Макет панели администратора (вкладка «Жанры»)}
    \label{fig:admin_genres}
\end{figure}